\documentclass[10pt,a4paper]{amsart}
\usepackage[margin=19mm]{geometry}
\usepackage{amsmath,amssymb,mathtools}
\usepackage[scr=rsfs,cal=euler]{mathalfa}
\usepackage{xfrac}
\usepackage[unicode,hidelinks,bookmarks=false]{hyperref}
\newcommand{\ie}{i.e.\ }
\newcommand{\cf}{cf.\ }
\newcommand{\resp}{respectively\ }
\newcommand{\Subsection}{\subsection}
\providecommand{\nobreakdash}{\nobreak\hbox{-}}
\setlength{\emergencystretch}{2em}
\multlinegap0pt
\usepackage{xcolor}
\usepackage{amssymb}
\usepackage{mathtools}
\usepackage[shortlabels]{enumitem}
\usepackage{tikz}
\usepackage{csquotes}
\newtheorem{proposition}{Proposition}
\newtheorem{theorem}{Theorem}
\newcommand{\cJ}{\mathcal{J}}
\newcommand{\cM}{\mathcal{M}}
\title{A vanishing theorem in Siefring's intersection theory}
\author{Naageswaran Manikandan}
\date{Source: arXiv:2412.11897v2 (CC BY 4.0)}
\begin{document}
\maketitle

\noindent\emph{Source and licence.} A vanishing theorem in Siefring's intersection theory. Version arXiv:2412.11897v2 (CC BY 4.0), distributed by arXiv under the Creative Commons Attribution 4.0 International licence, http://creativecommons.org/licenses/by/4.0/, as recorded in the arXiv licence metadata for this exact version and shown on the abstract page https://arxiv.org/abs/2412.11897v2. Full licence notice: \texttt{licenses/arxiv-2412.11897-CC-BY-4.0.txt}.\par\smallskip

\noindent\textit{Editorial scope.} Counted unit: the article's theorem delta\_rel=0 (source0.tex lines 1034--1037). The article's complete original proof of it is delivered with it (source0.tex lines 1039--1047, one \texttt{proof} environment of 9 lines). No mathematics is written, repaired or re-derived: the statement and the proof are the article's own lines, in the article's own notation, and no retained line is substituted or re-flowed.  Retained uncounted as the source-local context the proof needs: lines 787--790.  Nothing was omitted from any retained span, so the package carries no omission record.  No retained line contains a citation, so the wrapper carries no bibliography and no bibliography entry.  No retained line contains a figure environment, an image-inclusion command or drawing code, so no image is delivered and the package declares no asset dependency.  Editorial formatting only, in addition: an AMS \texttt{amsart} preamble that restates the article's own declarations and macros used by the retained lines, namely the environments \texttt{proposition}, \texttt{theorem} on separate counters, together with the standard packages those lines need.  The retained environments print in this document as Proposition 1, Theorem 1 in order of appearance, whereas the article prints the same items with the numbering of its own full document.  The complete native source of the article as distributed in the arXiv e-print for this exact version is bundled unchanged as \texttt{sources/170207/source0.tex}.
\begin{proposition}   
\label{regularity of evaluation map}
For any $z\in \Gamma_+$ and $\lambda = \sigma_{\max}^-(\gamma_z^{m_z})$, the evaluation map $e_{\lambda,z}\colon \cM(\widehat{W},\cJ^{\epsilon}_O) \rightarrow E_{\lambda,z}$ is submersion.
\end{proposition}

\begin{theorem} \label{delta_rel=0}
For any pair $z,w\in \Gamma_+$, there exist a comeagre subset of $\cJ_O$ such that for every $J$ in that comeagre subset, the subset of pseudoholomorphic curves in $\cM^*(\widehat{W},J)$ with $\delta_{\infty}([\tilde{u}; z], [\tilde{u};w])=0$ contains an open dense subset.

\end{theorem}

\begin{proof}
The proof can be divided into two parts. Let's start with a simpler case where 
$$m_w\cdot \sigma^-_{\max}(\gamma_z^{m_z}) \neq m_z \cdot \sigma^-_{\max}(\gamma_w^{m_w}).$$
In this case, the submanifold $\cM^*(\widehat{W},J)\setminus \left(\cM^z_{\lambda_2}(\widehat{W},J) \cup \cM^w_{\mu_2}(\widehat{W},J)\right)$ is an open dense subset of the moduli space $\cM^*(\widehat{W},J)$ and every pseudoholomorphic curve in it has $\delta_{\infty}([\tilde{u}; z], [\tilde{u};w])=0$.\\

Now, let's consider the case where $m_w\cdot \sigma^-_{\max}(\gamma_z^{m_z}) = m_z \cdot \sigma^-_{\max}(\gamma_w^{m_w})$. In this case, the submanifold $\cM^*(\widehat{W},J)\setminus \left(\cM^z_{\lambda_2}(\widehat{W},J) \cup \cM^w_{\mu_2}(\widehat{W},J)\right)$ might still contain curves with $\delta_{\infty}([\tilde{u}; z], [\tilde{u};w])\neq0$. But, the pseudoholomorphic curves in the open subset 
$$\left( \cM^*(\widehat{W},J)\setminus \left(\cM^z_{\lambda_2}(\widehat{W},J) \cup \cM^w_{\mu_2}(\widehat{W},J)\right) \right) \setminus (\bigcup\limits_{i=1}^{m_z\cdot m_w}h_i^{-1}(0))$$
have $\delta_{\infty}([\tilde{u}; z], [\tilde{u};w])=0$. Thus it would be sufficient to prove that the maps $h_{i}$'s are regular as this would mean the submanifolds $h_{i}^{-1}(0)$ have positive codimension. This regularity is proved the same way as Theorem~\ref{regularity of evaluation map} with minor modifications to the proof.
\end{proof}

\end{document}
